\subsubsection*{Solution}

{\color{blue}
We treat u1 and u2 as unknown parameters (instead of d=u1-u2 is fixed) and regard
\[
\theta=\frac13u1+\frac23u2
\]
as the parameter of interest.
}

In the weak-source limit, conditioning on detecting one photon gives the incoherent one-photon mixture used in quantum-superresolution theory \href{\detokenize{https://journals.aps.org/prx/abstract/10.1103/PhysRevX.6.031033}}{Tsang, Nair, and Lu}.

\subparagraph*{1. One-photon state}

Let

\[
|\psi_j\rangle = \int_{-\infty}^{\infty}dx\,\psi(x-u_j)|x\rangle, \qquad j=1,2,
\]

with $\langle\psi_j|\psi_j\rangle=1$. Since $\epsilon_0\to0$, the image-plane state in one temporal mode is, to first order,

\[
\rho\simeq (1-2\epsilon_0)|0\rangle\langle0| + \epsilon_0 \left( |\psi_1\rangle\langle\psi_1| + |\psi_2\rangle\langle\psi_2| \right).
\]

The vacuum probability is independent of the source positions. Conditional on measuring one photon, the normalized state is therefore

\[
\rho_1 = \frac{1}{2} \left( |\psi_1\rangle\langle\psi_1| + |\psi_2\rangle\langle\psi_2| \right).
\]

Thus the QFI of $\rho_1$ is the QFI per measured photon.

%the maths below is done by myself with the help of claude and latex code is generated by gpt-5.6 sol
{\color{blue}
\subparagraph*{2. Overlap identities}

Define the overlap directly in terms of the two source positions,
\[
\delta\equiv\langle\psi_1|\psi_2\rangle
=\int_{-\infty}^{\infty}dx\,\psi(x-u_1)\psi(x-u_2).
\]
For a real, normalized, sufficiently smooth PSF that vanishes at infinity,
\[
\langle\psi_i'|\psi_i'\rangle=\Delta k^2,
\qquad
\langle\psi_i|\psi_i'\rangle=0,
\]
and integration by parts gives
\[
\langle\psi_1'|\psi_2\rangle=\gamma,
\qquad
\langle\psi_1|\psi_2'\rangle=-\gamma,
\]
where, after shifting the integration variable,
\[
\gamma=\int_{-\infty}^{\infty}dx\,\psi'(x)\psi(x-u_2+u_1).
\]
Thus $\gamma$ measures how the overlap changes under a relative displacement of the two sources.

\subparagraph*{3. Two-parameter QFI matrix}

The two nonzero eigenstates and eigenvalues of $\rho_1$ are
\[
|e_\pm\rangle=
\frac{|\psi_1\rangle\pm|\psi_2\rangle}
{\sqrt{2(1\pm\delta)}},
\qquad
\lambda_\pm=\frac{1\pm\delta}{2}.
\]
Using the SLD quantum Fisher information matrix
\[
\mathcal Q_{ij}
=
2\sum_{m,n:\,\lambda_m+\lambda_n>0}
\frac{
\langle e_m|\partial_i\rho_1|e_n\rangle
\langle e_n|\partial_j\rho_1|e_m\rangle
}
{\lambda_m+\lambda_n},
\qquad i,j\in\{u_1,u_2\},
\]
and the overlap identities above, one obtains
\[
\boxed{
\mathcal Q_{\boldsymbol u}
=
\begin{pmatrix}
2\Delta k^2-\gamma^2 & -\gamma^2\\[3pt]
-\gamma^2 & 2\Delta k^2-\gamma^2
\end{pmatrix}},
\qquad
\boldsymbol u=(u_1,u_2)^\mathsf T.
\]

As a consistency check, a common displacement of both sources has QFI
\[
(1,1)\,\mathcal Q_{\boldsymbol u}\,(1,1)^\mathsf T
=4(\Delta k^2-\gamma^2),
\]
which is the fixed-separation translation result quoted in the proposed solution.

\subparagraph*{4. QFI for the weighted position $\theta$}

Write
\[
\theta=\boldsymbol v^\mathsf T\boldsymbol u,
\qquad
\boldsymbol v=
\begin{pmatrix}
1/3\\[2pt]
2/3
\end{pmatrix}.
\]
When both $u_1$ and $u_2$ are unknown, the multiparameter SLD Cram\'er--Rao bound for estimating this linear combination is
\[
\operatorname{Var}(\hat\theta)
\ge
\boldsymbol v^\mathsf T
\mathcal Q_{\boldsymbol u}^{-1}
\boldsymbol v.
\]
It is therefore natural to define the nuisance-adjusted information for $\theta$ by
\[
\mathcal Q_{\theta,\mathrm{eff}}
\equiv
\left(
\boldsymbol v^\mathsf T
\mathcal Q_{\boldsymbol u}^{-1}
\boldsymbol v
\right)^{-1}.
\]

The determinant is
\[
\det \mathcal Q_{\boldsymbol u}
=4\Delta k^2(\Delta k^2-\gamma^2),
\]
so
\[
\mathcal Q_{\boldsymbol u}^{-1}
=
\frac{1}{4\Delta k^2(\Delta k^2-\gamma^2)}
\begin{pmatrix}
2\Delta k^2-\gamma^2 & \gamma^2\\[3pt]
\gamma^2 & 2\Delta k^2-\gamma^2
\end{pmatrix}.
\]
Hence
\[
\begin{aligned}
\boldsymbol v^\mathsf T
\mathcal Q_{\boldsymbol u}^{-1}
\boldsymbol v
&=
\frac{1}{4\Delta k^2(\Delta k^2-\gamma^2)}
\left[
\frac{5}{9}(2\Delta k^2-\gamma^2)
+\frac{4}{9}\gamma^2
\right]\\[3pt]
&=
\frac{10\Delta k^2-\gamma^2}
{36\Delta k^2(\Delta k^2-\gamma^2)}.
\end{aligned}
\]
Therefore the QFI per measured photon, with the second source-position degree of freedom treated as a nuisance parameter, is
\[
\boxed{
\mathcal Q_{\theta,\mathrm{eff}}^{(1)}
=
\frac{36\Delta k^2(\Delta k^2-\gamma^2)}
{10\Delta k^2-\gamma^2}
}.
\]

For $N$ independent detected photons,
\[
\mathcal Q_{\theta,\mathrm{eff}}^{(N)}
=N\mathcal Q_{\theta,\mathrm{eff}}^{(1)},
\]
and the corresponding SLD Cram\'er--Rao bound is
\[
\operatorname{Var}(\hat\theta)
\ge
\frac{10\Delta k^2-\gamma^2}
{36N\Delta k^2(\Delta k^2-\gamma^2)}.
\]

}